\documentclass[10pt]{article}

\def\Type{LessonCaelan}
\def\TypeNum{Lesson 18}
\def\TypeTitle{Fundamental Theorem of Calculus \\Key }
\def\Semester{Spring 2026} %Not Used
\def\Course{MATH 1190 \& 1210}
\def\Targets{  \bf\color{Fuchsia}I3, \bf\color{Fuchsia}I4 }

\usepackage{preambleDJ}

%\usepackage{ragged2e}


%%%%%%%%%%%% BEGIN DOCUMENT %%%%%%%%%%%%%%%
%%%%%%%%%%%% BEGIN DOCUMENT %%%%%%%%%%%%%%%
%%%%%%%%%%%% BEGIN DOCUMENT %%%%%%%%%%%%%%%
%%%%%%%%%%%% BEGIN DOCUMENT %%%%%%%%%%%%%%%
%%%%%%%%%%%% BEGIN DOCUMENT %%%%%%%%%%%%%%%




\begin{document}

\pagestyle{otherpages}
\thispagestyle{firstpage}
\vs{.85}
{\bf\color{blue}Warm-up Exercise 1.}\ltrm{\bf\color{Fuchsia}I3} Find the following indefinite integrals.
\def\vd{.3}
\tasksA{2}{
\task $\ds \int e^x \,dx=$
\sol{ $\ds \int e^x \,dx=e^x+C$}
\task $\dint 3^x \,dx=$ 
\sol{ $\ds \int 3^x \,dx=\frac{3^x}{\ln(3)}+C$}
\vs{\vd}
\task  $\ds \int \frac{1}{x^2} \,dx =$ 
\sol{ \\\\
$\ds \int \frac{1}{x^2} \,dx=\dint x^{-2} \,dx=\frac{x^{-1}}{-1}+C=-\frac{1}{x}+C$}
\task $\ds \int \frac{1}{x} \,dx =$
\sol{$\ds \int \frac{1}{x} \,dx=\ln|x|+C$}
 \vs{\vd}
\task $\ds \int 2^5 \,dx =$ 
\sol{$2^5$ is a constant. \\
So $\ds \int 2^5 \,dx =2^5 \cdot x+C$ }
\task $\ds \int \log(2) \,dx =$ 
\sol{
$\log(2) $ is a constant.\\
So $\ds \int \log(2) \,dx = \log(2) \cdot x+C$
}\vs{\vd}
\task $\dint \frac23x^4 \,dx =$
\sol{\\\\
$\ds \int \frac23x^4 \,dx =\frac{2}{3}\cdot \frac{x^5}{5}+C=\frac{2}{15}x^5+C$
}
\task $\ds \int \sqrt{x^3} \,dx =$
\sol{\\\\
$\ds \int \sqrt{x^3} \,dx =\dint x^{3/2} \,dx =\frac{x^{5/2}}{5/2}+C=\frac{2}{5}x^{5/2}+C$
} \vs{\vd}
}
\newpage
\blue{\bf Mini-lecture:  The Fundamental Theorem of Calculus and its hypotheses}

%maybe start off with several simple ones - like 4x^3, 1/x, etc. - and then build up to larger ones - sums of these kinds of functions - then incorporate simplification

\begin{minipage}[t]{0.6\textwidth}
    \important{Fundamental Theorem of Calculus (FTC)}
    {
    If $f$ is continuous on $[a,b]$, then
    \vspace{-.1in}
    
    \[
    \int_a^b f(x)\ dx = F(x)\Big|_{a}^{b}
    \]
    
      where $F$ is any antiderivative of $f$ and\\ $F(x)\Big|_{a}^{b} = \Delta F = F(b)-F(a)$.
    }
\end{minipage}
%
\hfill
%
\begin{minipage}[t]{0.375\textwidth}
\,\\[-0.4in]
    \begin{center}
        \begin{tikzpicture}
          \begin{axis}[
        %    axis y line = left,
        %    axis x line = bottom,
        x label style={anchor=west},
            xlabel = {$x$},
            y label style={anchor=south},
            ylabel={$y$},
            %xtick       = {0.5,4.45},
            %xticklabels = { , },
            ytick       = {2.5,4.5,6.5,8.5},
            yticklabels = {$2$,$4$,$6$,$8$},
            samples     = 160,
            domain      = -2.5:2.5,
            restrict y to domain=-0.5:10,
            xmin = -2.5, xmax = 2.5,
            ymin = -0.5, ymax = 10,
            width = 3in, height=2in
          ]
          % piece 1
            %  \addplot[name path=f1, domain=-2.5:-0.01,  white, ultra thick,mark=none] {1/x^2+.5};
                           \addplot[name path=f1, domain=-2.5:-0.3, blue, ultra thick, <->] {1/x^2+.5};
  \path[name path=axis] (axis cs:-2.5,0) -- (axis cs:-0.3,0);
     \addplot [
        thick,
        color=blue,
        fill=blue, 
        fill opacity=0.1
    ]
    fill between[
        of=f1 and axis,
        soft clip={domain=-2.5:-0.3},
   ];
 \draw[color = blue, opacity =0.1, fill = blue, fill opacity=0.1] (-.3,0) rectangle (0,11.6);

          % piece 2
\addplot[name path=f2, domain=0.3:1, blue, ultra thick, <-] {1/x^2+.5};
  \path[name path=axis]  (axis cs:0.3,0)--(axis cs:1,0) ;
     \addplot [
        thick,
        color=blue,
        fill=blue, 
        fill opacity=0.1
    ]
    fill between[
        of=f2 and axis,
        soft clip={domain=0.3:1},
   ];
 \draw[color = blue, opacity =0.1, fill = blue, fill opacity=0.1] (.3,0) rectangle (0,11.6);

 % piece 3
\addplot[name path=f3, domain=1:2, red, ultra thick, mark=none] {1/x^2+.5};
  \path[name path=axis]  (axis cs:1,0)--(axis cs:2,0) ;
     \addplot [
        thick,
        color=red,
        fill=red, 
        fill opacity=0.1
    ]
    fill between[
        of=f3 and axis,
        soft clip={domain=1:2},
   ];
   
   \addplot[name path=f2, domain=2:2.5, blue, ultra thick, ->] {1/x^2+.5};
          \node at (1.5,6) {$f(x)=\dfrac1{x^2}$};
          \end{axis}
        \end{tikzpicture}
    \end{center}
\end{minipage}\\
%
\blue{\bf Example/Exercise}\ltrm{\bf\color{Fuchsia}I4} For each of the definite integrals below, a) Decide if $\ds \int_a^b \frac{1}{x^2}\,dx$ should be positive or negative. b) Calculate each definite integral. c) Discuss if your answer to a) and b) are consistent.
%\vs{-.2}
\begin{minipage}[t]{0.49\textwidth}
  \enumA{
  \item Based on graph, should $\dint_{1}^{2} \frac{1}{x^2}\,dx$ be positive or negative?\\
  \sol{Positive because on interval $[1,2]$ all area is above horizontal axis. See red shaded area above.}
  \item  $\dint_{1}^{2} \frac{1}{x^2}\,dx=$
  \sol{
  \al{
  \int_{1}^{2} \frac{1}{x^2}\,dx&=\int_{1}^{2} x^{-2}\,dx\\
  &=\frac{x^{-1}}{-1} \Big|_1^2\\
    &=\frac{-1}{x}\Big|_1^2\\
    &=\frac{-1}{2}-\frac{-1}{1}\\
    &=\frac{1}{2}
  }
  }
  \vs{-.3}
  \item Is your answer to a) consistent with b)? If not, what about the FTC might help us understand why not?\\
  \sol{Yes, a) and b) seem to be consistent with each other since both answers are positive.}
}
\end{minipage}
%
\vline\hfill
%
\begin{minipage}[t]{0.49\textwidth}
  \enumA{
  \item Based on graph, should $\dint_{-2}^{1} \frac{1}{x^2}\,dx$ be positive or negative?
    \sol{Positive because on interval $[-2,1]$ all area is above horizontal axis. See blue shaded area above.}
  \item $\dint_{-2}^{1} \frac{1}{x^2}\,dx=$
  \sol{
  \al{
  \int_{-2}^{1} \frac{1}{x^2}\,dx&=\int_{-2}^{1} x^{-2}\,dx\\
  &=\frac{x^{-1}}{-1} \Big|_{-2}^1\\
    &=\frac{-1}{x}\Big|_{-2}^1\\
    &=\frac{-1}{1}-\frac{-1}{-2}\\
    &=-\frac{3}{2}
 }
  }
  \vs{-.3}
  \item Is your answer to a) consistent with b)? If not, what about the FTC might help us understand why not?\\
  \sol{No, area should be positive from a) but in b) we found area to be negative.  The issue is that   $f(x)=\frac{1}{x^2}$ is {\em \bf not} continuous at $x=0$ due to division by $0$. Further, $x=0$ is in the domain $[-2,1]$  and thus violates FTC hypothesis that $f(x)$ must be continuous on $[-2,1]$. So {\em FTC does not apply} and we need to use other techniques to evaluate $\dint_{-2}^{1} \frac{1}{x^2}\,dx$.}
}
\end{minipage}
\newpage
\exercise \ltrm{\bf\color{Fuchsia}I4}For each of the following definite integrals, decide if the Fundamental Theorem applies.  If it does, evaluate the integral.\\

\tasksA{2}{
\task $\displaystyle\int_0^2 2x \,dx$ \qquad FTC apply? \sol{Yes! $f(x)=2x$ is continuous everywhere which means continuous on $[0,2]$. Practically speaking, we don't have a division by $0$, square root of $(-)$ values, log of $0$ or $(-)$ values, etc.}
\task $\displaystyle\int_{-2}^{-1} \sqrt{x} \,dx$ \qquad FTC apply? \sol{No! $f(x)=\sqrt{x}$ is not continuous for negative numbers which occur on domain $[-2,-1]$.   Thus we don't meet hypothesis of FTC that $f(x)=\sqrt{x}$ be continuous on $[-2,-1]$ and we cannot use it.}
}
\tasksI{2}{
\task If ``yes", \\$\displaystyle\int_0^2 2x \,dx$=
\sol{\al{
\int_0^2 2x \,dx&=x^2\Big|_0^2\\
&=2^2-0^2\\
&=4
}
}
\task If ``yes", \\$\displaystyle\int_{-2}^{-1} \sqrt{x} \,dx=$
}


%%%%%%%%%    
%\newpage
%%%%%%%%%

\example \ltrm{\bf\color{Fuchsia}I4}Evaluate the following integrals: {\em Hint:} It might be helpful to a bit of algebra first so that you can use one of our derivative rules.

\begin{itemize}

\item[(a)] $\ds \int_1^2 \left( \frac{2}{\sqrt[5]{t^3}} + 3^t - \frac12 \right)\,dt$
\sol{ We note that $\ds f(t)=\dfrac{2}{\sqrt[5]{t^3}} + 3^t - \frac12$ is discontinuous at $t=0$ since we have a division by $0$.\\ \\ But $t=0$ is not in the domain $[1,2]$. Thus FTC applies since $f(t)$ is continuous on $[1,2]$.  
\al{
\int_1^2 \left( \dfrac{2}{t^{3/5}} + 3^t - \frac12 \right)\,dt&=\int_1^2 \left( 2t^{-3/5} + 3^t - \frac12 \right)\,dt\\
&=2\frac{t^{2/5}}{2/5}+\frac{3^t}{\ln(3)}-\frac{1}{2} t \Big|_1^2\\
&=\left( 2\frac{2^{2/5}}{2/5}+\frac{3^2}{\ln(3)}-\frac{1}{2}\cdot  2 \right) -\left( 2\frac{1^{2/5}}{2/5}+\frac{3^1}{\ln(3)}-\frac{1}{2} \cdot 1  \right)
}
{\em Note}: It's not necessary to simplify!
}

\item[(b)] $\dint_2^4 \dfrac{z^2-z-1}{2z}\,dz$
\sol{
We have a division by $0$ (i.e. discontinuity) when $z=0$ for $f(z)= \dfrac{z^2-z-1}{2z}$. But $z=0$ is not in the domain $[2,4]$ so FTC applies.
\al{
\int_2^4 \frac{z^2-z-1}{2z}\,dz&=\int_2^4 \frac{z^2}{2z}-\frac{z}{2z}-\frac{1}{2z}\,dz\\
&=\int_2^4 \frac{z}{2}-\frac{1}{2}-\frac{1}{2}\cdot \frac{1}{z}\,dz\\
&=\frac{1}{2}\cdot \frac{z^2}{2}-\frac{1}{2}z-\frac{1}{2}\ln|z| \Big|_2^4\\
&= \frac{z^2}{4}-\frac{z}{2}-\frac{1}{2}\ln|z| \Big|_2^4\\
&=\left(  \frac{4^2}{4}-\frac{4}{2}-\frac{1}{2}\ln|4|\right) - \left( \frac{2^2}{4}-\frac{2}{2}-\frac{1}{2}\ln|2| \right)
}
}

\end{itemize}
\newpage
\exercise \ltrm{\bf\color{Fuchsia}I4}Evaluate the following integrals:

\begin{itemize}

\item[(a)] $\ds \int_{-2}^{-1} \left( \frac{4}{\sqrt[3]{u}} + e^u + 2u^3 \right) \, du$

\sol{ $\ds f(u)= \frac{4}{\sqrt[3]{u}} + e^u + 2u^3$ is discontinuous at $u=0$ (division by $0$). But $u=0$ is not in the domain $[-2,-1]$ so FTC applies.
\al{
\int_{-2}^{-1} \left( \frac{4}{\sqrt[3]{u}} + e^u + 2u^3 \right) \, du&=\int_{-2}^{-1} \left( \frac{4}{u^{1/3}} + e^u + 2u^3 \right) \, du \\
&=\int_{-2}^{-1} \left( 4u^{-1/3} + e^u + 2u^3 \right) \, du\\
&=4\frac{u^{2/3}}{2/3} +e^u+2\cdot \frac{u^4}{4} \Big|_{-2}^{-1}\\
&=6u^{2/3}+e^u +\frac{u^4}{2}   \Big|_{-2}^{-1}\\
&=\left( 6(-1)^{2/3}+e^{-1} +\frac{(-1)^4}{2}   \right) -\left( 6(-2)^{2/3}+e^{-2} +\frac{(-2)^4}{2}   \right)
}

}
\item[(b)] $\displaystyle\int_{1}^2 \frac{1+\sqrt[4]{y}+\sqrt{y}+y}{y} \, dy$
\sol{
At $y=0$, $\ds f(y)=\frac{1+\sqrt[4]{y}+\sqrt{y}+y}{y} $ is discontinuous (division by $0$).  But since $y=0$ is not in the domain $[1,2]$, we can use the FTC.
\al{
\int_{1}^2 \frac{1+\sqrt[4]{y}+\sqrt{y}+y}{y} \, dy&=\int_{1}^2 \frac{1}{y}+\frac{y^{1/4}}{y}+\frac{y^{1/2}}{y}+\frac{y}{y} \, dy\\
&=\int_{1}^2 \frac{1}{y}+y^{-3/4}+y^{-1/2}+1 \, dy\\
&=\ln|y| +\frac{y^{1/4}}{1/4}+\frac{y^{1/2}}{1/2}+y \Big|_1^2\\
&= \ln|y|+4 y^{1/4}+2 y^{1/2}+y \Big|_1^2\\
&= \left(   \ln|2|+4 \cdot 2^{1/4}+2\cdot  2^{1/2}+2   \right) -\left(  \ln|1|+4  \cdot1^{1/4}+2 \cdot1^{1/2}+1  \right)
}
}

\end{itemize}

%%%%%%%%%    
\newpage%
%%%%%%%%%

{\bf\color{blue}Extra Practice 1.} Evaluate the following integrals:

\begin{itemize}

\item[(a)]\ltrm{\bf\color{Fuchsia}I3} $\ds \int \left(x^2+x+1\right)\,dx$ 
\sol{ 
\al{
 \int \left(x^2+x+1\right)\,dx&=\frac{x^3}{3}+\frac{x^2}{2}+x+C
}
Note that the ``$+C$" is required because indefinite integral is asking for ``the most general anti-derivative."
}

\item[(b)] \ltrm{\bf\color{Fuchsia}I4}$\dint_0^1 \left(x^2+x+1\right)\,dx$
\sol{
$f(x)=x^2+x+1$ is continuous for all values of $x$ (i.e. not division by $0$, square root of $(-)$ values, no logarithms of $0$ or  $(-)$ values, etc.). So we can utilize the FTC
\al{
 \int \left(x^2+x+1\right)\,dx&=\frac{x^3}{3}+\frac{x^2}{2}+x\Big|_0^1\\
 &=\left(\frac{1^3}{3}+\frac{1^2}{2}+1  \right)-\left(\frac{0^3}{3}+\frac{0^2}{2}+0  \right)
\intertext{ It's ok to stop after previous line without further simplification.}
 &=\left(\frac{1}{3}+\frac{1}{2}+1\right)-0\\
 &=\frac{2}{6}+\frac{3}{6}+\frac{6}{6}\\
 &=\frac{11}{6}
}
}
\item[(c)] \ltrm{\bf\color{Fuchsia}I4}$\ds \int_{1/8}^1 \left( \frac{x^2+x+1}{\sqrt[3]{x}} \right)\,dx$
\sol{
$\ds f(x)=  \dfrac{x^2+x+1}{\sqrt[3]{x}}$ is discontinuous at $x=0$ (division by $0$). But $x=0$ is not in domain $[1/8,1]$ so FTC applies.
\al{
\int_{1/8}^1 \left( \frac{x^2+x+1}{\sqrt[3]{x}} \right)\,dx&=\int_{1/8}^1 \left( \frac{x^2}{x^{1/3}}+\frac{x}{x^{1/3}}+\frac{1}{x^{1/3}} \right)\,dx\\
&=\int_{1/8}^1 x^{5/3}+x^{2/3}+x^{-1/3}\,dx\\
&=\frac{x^{8/3}}{8/3}+\frac{x^{5/3}}{5/3}+\frac{x^{2/3}}{2/3} \Bigg|_{1/8}^1\\
&=\frac{3}{8}x^{8/3}+\frac{3}{5}x^{5/3}+\frac{3}{2}x^{2/3} \Bigg|_{1/8}^1\\
&=\left( \frac{3}{8}1^{8/3}+\frac{3}{5}1^{5/3}+\frac{3}{2}1^{2/3}  \right) -\left( \frac{3}{8}\left(\frac{1}{8}\right)^{8/3}+\frac{3}{5}\left(\frac{1}{8}\right)^{5/3}+\frac{3}{2}\left(\frac{1}{8}\right)^{2/3} \right)  
%\intertext{ It's ok to stop after previous line without further simplification.}
%&=\left(\frac{3}{8}+\frac{3}{5}+\frac{3}{2} \right) -\left( \frac{3}{8}\cdot\frac{1}{2^8}+\frac{3}{5}\cdot \frac{1}{2^5}+\frac{3}{2} \cdot \frac{1}{2^2} \right)
}
}
\end{itemize}

\newpage
{\bf\color{blue}Extra Practice 2.} \ltrm{\bf\color{Fuchsia}I4}The graph of $f(x)=x^3$ is given below.\\
\mpageT{.35}{.65}{    
\begin{center}
        \begin{tikzpicture}
          \begin{axis}[
       xmin=-2.5,        % Minimum x-value
            xmax=2.5,         % Maximum x-value
            ymin=-8,         % Minimum y-value
            ymax=8,          % Maximum y-value
            axis lines=middle, % Draw the x and y axes in the middle
            xlabel={$x$},     % Label for the x-axis
            ylabel={$y$},     % Label for the y-axis
            xtick       = {-2,-1,1,2},
            ytick       = {-8,-4,4,8},
            %yticklabels = {$2$,$4$,$6$,$8$},
            axis line style={thick, <->}, % Add arrows at the end of the axes
            width = 2in, height=2in,
            %restrict y to domain=-9:9
          ]
          %\addplot[domain=-2.5:2.5, samples=400, blue, ultra thick] {x^3};
              \addplot[name path=f1, domain=-2.5:2.5, blue, ultra thick, <->] {x^3};
  \path[name path=axis] (axis cs:-2.5,0) -- (axis cs:2.5,0);
     \addplot [
        thick,
        color=blue,
        fill=blue, 
        fill opacity=0.1
    ]
    fill between[
        of=f1 and axis,
        soft clip={domain=-1:2},
   ];
          \end{axis}
        \end{tikzpicture}
    \end{center}
}{
%
\begin{itemize}
\item[(a)] Find the {\bf net area} between the graph of $f(x) = x^3$ and the portion of the horizontal axis from $x=-1$ to $x=2$.
\sol{
We want $\ds \int_{-1}^2 x^3\;dx$. Since $f(x)=x^3$ is continuous everywhere, FTC applies. 
\al{
\int_{-1}^2 x^3\;dx&=\frac{x^4}{4}\Bigg|_{-1}^{2}\\
&=\frac{2^4}{4}-\frac{(-1)^4}{4}\\
&=4-\frac{1}{4}\\
&=\frac{15}{4}
}
}
\end{itemize}
}
\begin{itemize}
\item[(b)] Find the {\bf area} between the graph of $f(x)=x^3$ and the portion of the horizontal axis from $x=-1$ to $x=2$.
\sol{
We're asked for {\bf area} which assumes a positive value. Since the area on $[-1,0]$ is negative, and area on $[0,2]$ is positive, we break apart our integral into two parts and force the negative part to be positive. Here's one way to do this:
\al{
\Bigg| \int_{-1}^0 x^3\;dx \Bigg| +\int_{0}^2 x^3\;dx&=\Bigg[ \Bigg| \frac{x^4}{4} \Big|_{-1}^0 \Bigg|  \Bigg]\,+\, \Bigg[\frac{x^4}{4}\Big|_0^2\Bigg]\\
&=\Bigg[ \Bigg| \frac{0^4}{4}-\frac{(-1)^4}{4}  \Bigg|  \Bigg]\,+\, \Bigg[\frac{2^4}{4} - \frac{0^4}{4}\Bigg]\\
&=\Bigg[ \frac{1}{4}  \Bigg]+\Bigg[ 4  \Bigg]\\
&=\frac{17}{4}
}
}
\end{itemize}

\end{document}
